% 方法段骨架。数字不要手写，用 \input{tabs/table_car.tex}。
% 需要 booktabs 与 threeparttable。

\section{实证设计}

我们以市场模型度量事件窗内的累计异常收益。
对每家公司，在事件日前的估计窗内拟合
\( R_{it} = \alpha_i + \beta_i R_{mt} + \varepsilon_{it} \)，
事件窗内的异常收益为 \( AR_{it} = R_{it} - \hat\alpha_i - \hat\beta_i R_{mt} \)。
估计窗必须落在所有报告窗口之前，避免事件日本身进入 \(\beta\) 的估计。
收益使用简单收益率；行情为前复权收盘价。
非交易日的事件日顺延到其后十个日历日内的第一个交易日。

横截面 \(t\) 统计量假设事件互相独立，只作对照。
事件日分散时，首选 Boehmer, Musumeci and Poulsen (1991) 的标准化检验；
估计窗残差存在截面相关时，再用 Kolari and Pynnönen (2010) 的相关调整。
若半数以上事件落在同一天，首选等权组合的时间序列检验，因为这一天是一次共同冲击，不是 \(N\) 次独立冲击。
平均 CAR 上的星号来自首选检验，没有做多重检验校正。
窗口表里的 p BH 只对这些窗口的首选 p 值做 Benjamini--Hochberg 调整，不改星号。
秩检验只作厚尾对照。

行业指数和三因子由使用者提供，程序不下载、也不自行匹配。
时间安慰剂在伪事件日上重新拟合估计窗，不是把真实事件的异常收益平移。
安慰剂不是对照组。
若事件表标明盘后公告，对齐从记录日的下一个交易日开始，仍限制在十个日历日以内。
涨跌停标记按板块和公司名称推断带宽，不是交易所状态表，不因此删行。

回归是事件层累计异常收益对披露或公司特征的横截面关联，标准误按公司聚类。
它不是平均处理效应。
信号、合法性与制度同构只作为估计前的可区分预测，写在 \texttt{output/theory/candidate\_mechanisms.md}；
结果出来之后再按完整的系数模式取舍，包括不显著的列。

\input{tabs/table_car.tex}
\input{tabs/table_summary.tex}
\input{tabs/table_regression.tex}
\input{tabs/table_placebo.tex}

\begin{figure}[htbp]
\centering
\includegraphics[width=0.85\textwidth]{figs/fig_car_path.pdf}
\caption{事件路径上的平均累计异常收益。说明见 \texttt{figs/captions.tex}。}
\label{fig:car_path}
\end{figure}
